\documentclass[conference]{IEEEtran}
\usepackage{amsmath,url}
\newcommand{\toprule}{\hline}\newcommand{\midrule}{\hline}\newcommand{\bottomrule}{\hline}
\title{Decision-Driven Evidence Disclosure for Collaborative Security Correlation}
\author{\IEEEauthorblockN{Ananay Srivastava, Jitendra Singh, Manik Bhardwaj, Ashish Jain, Mukunj Singh}\IEEEauthorblockA{Department of Computer Science and Engineering, Sharda University, India}}
\begin{document}\maketitle
\begin{abstract}Collaborative detection needs corroboration, but sharing more telemetry increases exposure and cost. We present a finite, policy-controlled disclosure protocol that requests only predicates capable of changing an unresolved hypothesis. Each release is versioned, coarsened or refused by its owner, debited transactionally, and charged against a cumulative budget across retries and overlapping investigations. A durable correlation worker counts independent provenance groups and records supported, contradicted, or unresolved outcomes. In a controlled protocol study of 1,000 deterministic cases, decision-driven disclosure reduced no-fault policy units from 660 to 540 and payload bytes from 230,460 to 200,100 versus fixed richer sharing while retaining 20/20 true-positive cases in the reported no-fault slice. Fault injections expose the cost of refusal and budget exhaustion. These are protocol results, not claims of differential privacy or real-organization accuracy; CIC-IDS2017 replay was not completed because the official archive endpoint required registration.\end{abstract}
\section{Introduction}Security collaboration is constrained by both incomplete visibility and unwillingness to export raw logs. Existing selective-sharing systems establish the value of policies and budgets, but a fixed release plan can disclose information that cannot change the current decision. Our question is narrower: can a coordinator choose the next approved aggregate from the unresolved predicates, enforce cumulative accounting, and stop explicitly when evidence is insufficient?
\section{Contribution and scope}The contribution is an implementable protocol, not a new detector. Signals carry stable identifiers, schema versions, observation and receipt intervals, detector version, and provenance. A finite catalogue contains authentication-burst and traffic-burst aggregates; raw-log queries are outside the interface. Correlation state is durable and owned by one worker. We report synthetic organization partitions as protocol cases, not independent organizations.
\section{Protocol}Let $H$ be the current hypothesis set and $P$ the predicates that distinguish competing hypotheses. At round $r$, the coordinator selects $p^*=\arg\min_{p\in P_H}(c_p,\mathrm{id}_p)$, where $c_p$ is a deterministic policy cost and the identifier breaks ties. An owner returns approve, coarsen, or deny. Approval debits a release ledger in the same transaction that records the immutable response; a retry reuses the response by request identifier. Limits aggregate by investigation, organization, and policy epoch, so a fresh request cannot reset a budget. The protocol stops on sufficient support, contradiction, exhaustion, expiry, or refusal of all useful predicates. Outcomes remain distinct from analyst workflow status.
\section{Implementation}The reference deployment uses FastAPI, PostgreSQL, one correlation worker, organization agents with SQLite outboxes, and authenticated browser sessions. The database stores an inbox, incident state, evidence revisions, disclosure requests, release ledger, and notification outbox. Authentication keys are hashed and mapped to enrolled nodes. Dashboard delivery uses bounded queues and recovery cursors. The design minimizes raw export but does not guarantee privacy: aggregates can still leak membership, timing, or rare-event information.
\section{Evaluation}We froze eight cases and ten seeds. Five strategies were compared: local-only, fixed minimal, fixed richer, static escalation, and decision-driven disclosure. Four fault classes were injected: none, refusal, coarsening, and budget exhaustion.\begin{table}[t]\centering\caption{No-fault protocol slice (50 cases per strategy).}\begin{tabular}{lrrrr}\toprule Strategy&TP&FP&Units&Bytes\\\midrule Local-only&20&20&0&0\\Fixed minimal&20&20&120&48,180\\Fixed richer&20&0&660&230,460\\Static escalation&20&0&660&230,460\\Decision-driven&20&0&540&200,100\\\bottomrule\end{tabular}\end{table}The decision-driven policy used 300 refinements and two rounds on average in this slice; fixed richer used 360 refinements and 0.8 rounds. Under denial or exhausted budgets, unresolved outcomes increased, which is the intended safe failure mode. The implementation also passed 37 automated tests, frontend checks, migration and backup/restore checks, and a three-organization local demo.
\section{Limitations and threats to validity}The protocol harness uses synthetic partitions and a simplified detector; it does not establish field accuracy, high availability, or formal privacy. The official CIC-IDS2017 replay remains incomplete. Measurements were taken on one machine and should be repeated with controlled hardware. Published selective disclosure and privacy-budget methods are close prior art; novelty therefore rests on the precise predicate-selection, cumulative-accounting, and outcome semantics together with comparative evidence.
\section{Conclusion}Decision-driven disclosure makes sharing conditional on decision value and gives refusal and exhaustion explicit semantics. The current evidence supports a reproducible protocol result, not a universal privacy or acceptance claim.
\begin{thebibliography}{00}\bibitem{kum}Kum et al., ``Privacy-preserving collaborative intrusion detection,'' SOUPS, 2019.\bibitem{dunnett}Dunnett et al., ``Privacy-preserving cyber threat intelligence sharing,'' arXiv:2208.12031, 2022.\bibitem{cic}Canadian Institute for Cybersecurity, ``CIC-IDS2017,'' University of New Brunswick, 2017.\bibitem{misp}MISP Project, ``MISP threat sharing platform,'' 2024.\end{thebibliography}\end{document}
